\documentclass[12pt,a4paper]{article}

\usepackage[utf8]{inputenc}
\usepackage[T1]{fontenc}
\usepackage[ngerman]{babel}
\usepackage[a4paper,margin=2.5cm]{geometry}
\usepackage{enumitem}
\usepackage{parskip}

\title{Exposé }
\author{}
\date{}

\begin{document}
	
	\maketitle
	
	\section*{Arbeitstitel}
	
	\textbf{Workflow-basierte Vorabdisposition von Heizöl-Bestellungen mit BPMN und Camunda}
	
	\section{Aufgabenstellung und Zielsetzung}
	
	Ziel des Projekts ist die prototypische Entwicklung eines Systems zur Unterstützung der Vorabkoordination von Heizöl-Bestellungen.
	
	Ein Kunde soll über ein Web-Frontend eine Bestellung erfassen können. Dabei werden unter anderem folgende Informationen übermittelt:
	
	\begin{itemize}[leftmargin=1.5cm]
		\item Bestellmenge
		\item Lieferadresse
		\item Kommunikationskanal
		\item Name
		\item Telefonnummer
		\item E-Mail-Adresse
		\item optionaler Kommentar
	\end{itemize}
	
	Auf Basis definierter Regeln soll das System anschließend einen passenden Lieferanten sowie einen möglichen Liefertermin ermitteln.
	
	Dieser Terminvorschlag wird dem Kunden per E-Mail/SMS/WhatsApp übermittelt. Über einen Bestätigungs- oder Ablehnungslink kann der Kunde auf den Vorschlag reagieren.
	
	Bei Bestätigung wird der Auftrag an die Disposition weitergegeben und dort sichtbar gemacht. Gleichzeitig wird der Status des Auftrags im System aktualisiert.
	
	Bei Ablehnung kann ein alternativer Vorschlag erzeugt oder der Vorgang beendet werden.
	
	Erfolgt innerhalb eines definierten Zeitraums keine Rückmeldung, wird der Auftrag automatisch storniert.
	
	\section{Ist-Situation und Problemstellung im Unternehmenskontext}
	
	Die Projektidee basiert auf einem praktischen Anwendungsfall aus einem Unternehmen, in dem Heizöl-Bestellungen aktuell überwiegend manuell koordiniert werden.
	
	Die Terminfindung sowie die Abstimmung mit Kunden erfolgen derzeit nach Einschätzung des Unternehmens im Wesentlichen manuell. Ein automatisierter Vorprozess zur strukturierten Erfassung, Terminermittlung und Rückmeldung existiert derzeit nicht.
	
	Das Projekt soll deshalb als prototypische Lösung untersuchen, wie ein solcher Ablauf technisch unterstützt und in Form eines klar definierten Workflows abgebildet werden kann.
	
	Produktive Unternehmensdaten werden dabei nicht verwendet; die Umsetzung erfolgt ausschließlich mit synthetischen Testdaten.
	
	\section{Technische Einordnung / Infrastruktur}
	
	Die technische Umsetzung ist als eigenständiger Prototyp vorgesehen.
	
	Geplant ist die Verwendung folgender Technologien:
	
	\begin{itemize}[leftmargin=1.5cm]
		\item Spring Boot als Backend-Basis
		\item Camunda 7 zur Workflow-Ausführung
		\item PostgreSQL zur Datenhaltung
		\item einfaches Web-Frontend zur Benutzerinteraktion
	\end{itemize}
	
	Die Lösung soll zunächst unabhängig von bestehender Unternehmenssoftware entwickelt werden.
	
	Technologische Vorgaben seitens des Unternehmens werden aktuell noch abgestimmt.
	
	\section{Geplante methodische Vorgehensweise und Ergebnistyp}
	
	Das Projekt wird schrittweise umgesetzt:
	
	\begin{enumerate}[leftmargin=1.5cm]
		\item Analyse und fachliche Abgrenzung des Prozesses
		\item Modellierung des Geschäftsprozesses in BPMN
		\item Implementierung des Workflows mit Camunda
		\item Entwicklung des Backends für Prozesslogik und Statusverwaltung
		\item Entwicklung eines Frontends zur Eingabe und Rückmeldung
		\item Test typischer Prozessszenarien
	\end{enumerate}
	
	Das geplante Ergebnis ist ein prototypisches System, bestehend aus:
	
	\begin{itemize}[leftmargin=1.5cm]
		\item einem ausführbaren BPMN-Prozess
		\item einem Backend-Service
		\item einem Frontend für Kundeninteraktion
		\item einer einfachen Ansicht für die Disposition
	\end{itemize}
	
	\section{Vorläufige Aufgabenverteilung im Projektteam}
	
	\subsection*{Daniel Tempel}
	
	Technische Projektkoordination, Backend-Architektur, REST-Schnittstellen, Integration von Camunda und Backend, Kommunikation mit dem Unternehmen sowie fachliche Betreuung.
	
	\subsection*{Buse}
	
	BPMN-Modellierung, Prozesslogik, Sonderfälle, Prozessvariablen, workflow-nahe Services und Testszenarien.
	
	\subsection*{Arsenii}
	
	Frontend, Benutzerinteraktion, Kunden- und Disponentenansicht, Dokumentationspflege sowie Testunterstützung.
	
	Zusätzlich erfolgen wöchentliche Teamabstimmungen jeweils freitags sowie Aufgabenverwaltung über GitHub (Issues, Branches, Pull Requests).
	
\end{document}